\documentclass[../../main.tex]{subfiles}
\begin{document}

{\bf[解]}

簡単のため
\begin{align}
  t=a-\dfrac{1}{a} \label{eq:1}
\end{align}
とおく．

\begin{figure}[htb]
  \centering
  \begin{tikzpicture}[scale=1.5, >=stealth, every node/.style={font=\small}]

    % --- 1. 座標定義 ---
    \coordinate (O)     at (0, 0);
    \coordinate (Pt_p1) at (1.0, 0);   % a = 1, t = 0
    \coordinate (Pt_m1) at (-1.0, 0);  % a = -1, t = 0

    % --- 2. 漸近線 t = a (斜め45度) ---
    \draw[dashed, gray!55, thick] (-2.8, -2.8) -- (2.8, 2.8)
    node[above right=-1pt, font=\scriptsize, gray!90!black] {漸近線 $t = a$};

    % --- 3. 座標軸 ---
    \draw[->, thick, gray!80!black] (-3.2, 0) -- (3.3, 0) node[right] {$a$};
    \draw[->, thick, gray!80!black] (0, -3.2) -- (0, 3.3) node[above] {$t$};
    \node[below left=2pt] at (O) {$O$};

    % --- 4. 曲線 t = a - 1/a の描画 ---
    % 右側の分枝 (a > 0)
    \draw[line width=1.8pt, blue!85!black, domain=0.31:2.9, samples=80]
    plot (\x, {\x - 1.0/\x});
    \node[blue!90!black, font=\footnotesize\bfseries, right] at (2.6, {2.6 - 1.0/2.6})
    {$t = a - \frac{1}{a}$};

    % 左側の分枝 (a < 0)
    \draw[line width=1.8pt, blue!85!black, domain=-2.9:-0.31, samples=80]
    plot (\x, {\x - 1.0/\x});

    % --- 5. a 切片 a = ±1 (t = 0, t^2 が最小となる点) ---
    \fill[red!85!black] (Pt_p1) circle (1.5pt);
    \fill[red!85!black] (Pt_m1) circle (1.5pt);
    \node[below right=1pt, font=\footnotesize\bfseries, red!90!black] at (Pt_p1) {$1$};
    \node[above left=1pt, font=\footnotesize\bfseries, red!90!black]  at (Pt_m1) {$-1$};

  \end{tikzpicture}
  \caption{関数 $t = a - \dfrac{1}{a}$ のグラフ}
  \label{fig:t_a_minus_inv_a}
\end{figure}

題意の関数は
\begin{align}
  f(x) = (3x^2-4)(x-t) = 3x^3 -3tx^2-4x+4t \label{eq:2}
\end{align}
だから，一階微分は
\begin{align}
  f'(x)=9x^2-6tx-4 \label{eq:3}
\end{align}
である．$f'(0)<0$および$2$次係数が正だから$f'(x)=0$は$2$つの異なる実解$\alpha$，$\beta$($\alpha<\beta$)を持つ．
よって$f(x)$の増減表は以下のようになる．

\begin{table}[htb]
  \centering
  \caption{$f(x)$ の増減表}
  \label{tab:fx_monotonicity}
  \renewcommand{\arraystretch}{1.4}
  \begin{tabular}{c|ccccc}
    $x$     & $\dots$    & $\alpha$   & $\dots$    & $\beta$    & $\dots$    \\ \hline
    $f'(x)$ & $+$        & $0$        & $-$        & $0$        & $+$        \\ \hline
    $f(x)$  & $\nearrow$ & \text{極大} & $\searrow$ & \text{極小} & $\nearrow$
  \end{tabular}
\end{table}

従って，極大値と極小値の差$L$として，
\begin{align}
  L=&f(\alpha)-f(\beta) \\
  =&\int_{\beta}^{\alpha}f'(x)\,dx \\
  =&\int_{\alpha}^{\beta}9(x-\alpha)(\beta-x)\,dx \\
  =&\dfrac{3}{2}(\beta-\alpha)^3
\end{align}
である．故に$\beta-\alpha(>0)$の最小値を求めればよい．

ここで，$\alpha$，$\beta$は$f'(x)=0$の$2$解であるから，\cref{eq:3}より
\begin{align}
  \beta-\alpha=\dfrac{2\sqrt{9t^2+36}}{9}=\dfrac{2\sqrt{t^2+4}}{3}
\end{align}
である．従って$t^2$が最小の時$L$は最小である．
$t=a-\dfrac{1}{a}$の図から，$t^2$が最小となるのは$t^2=0$の時で，この時
\begin{align}
  a = \pm 1
\end{align}
である．$\cdots$(答)
\newpage
\end{document}
